\chapter{银牌档高频模板}

本章只收录区域赛中值得提前准备、但不适合临场重新推导的模板。每个模板都必须在复制后用本节样例和边界测试跑过。

\section{莫队算法}

\cardmeta{现场常用}{离线区间统计，移动左右端点的代价小}{排序 \(O((n+q)\sqrt n)\)}{不能处理在线询问；必须实现对称的 add/remove}

\begin{lstlisting}[style=cpp]
struct MoQuery { int l, r, id; };
int block_size;
vector<int> a, freq; // a 是离散化后的数组，freq[x] 是当前窗口中 x 的次数
int distinct_count;

void add_pos(int p) {
    // 把位置 p 纳入窗口；第一次出现时不同数个数 +1
    if (++freq[a[p]] == 1) ++distinct_count;
}
void remove_pos(int p) {
    // 移除位置 p；最后一次消失时不同数个数 -1
    if (--freq[a[p]] == 0) --distinct_count;
}
vector<int> mo_distinct(vector<MoQuery> qs, int n) {
    // 莫队把查询排序，使窗口左右端点移动总量尽量小
    block_size = max<int>(1, sqrt(n));
    sort(qs.begin(), qs.end(), [&](auto x, auto y) {
        int bx = x.l / block_size, by = y.l / block_size;
        if (bx != by) return bx < by;
        return (bx & 1) ? x.r > y.r : x.r < y.r;
    });
    freq.assign(n + 1, 0);
    distinct_count = 0;
    vector<int> ans(qs.size());
    int l = 0, r = -1; // 初始窗口为空，维护闭区间 [l,r]
    for (auto q : qs) {
        while (l > q.l) add_pos(--l);
        while (r < q.r) add_pos(++r);
        while (l < q.l) remove_pos(l++);
        while (r > q.r) remove_pos(r--);
        ans[q.id] = distinct_count;
    }
    return ans;
}
\end{lstlisting}

\begin{card}
\textbf{完整样例}\\
\begin{lstlisting}
5 3
1 2 1 3 2
0 2
1 4
0 4
\end{lstlisting}
输出：
\begin{lstlisting}
2
3
3
\end{lstlisting}
这里询问使用 0-based 闭区间；数组离散化后再维护 `freq`。
第四组边界测试使用空区间、单点区间和全部相同数组。
\end{card}

\section{可持久化线段树：区间第 \(k\) 小}

\cardmeta{现场常用}{历史版本、静态区间第 \(k\) 小}{单次查询 \(O(\log n)\)}{值域先离散化；每个版本只复制访问路径}

\begin{lstlisting}[style=cpp]
struct PNode { int l = 0, r = 0, sum = 0; };
vector<PNode> pst(1); // 0 号空节点，未走到的子树都指向它
vector<int> roots, values;

int update(int old, int L, int R, int pos) {
    // 新版本复制旧版本的路径，其余节点与旧版本共享
    int cur = pst.size();
    pst.push_back(pst[old]);
    ++pst[cur].sum; // 当前值域区间内出现次数 +1
    if (L != R) {
        int M = (L + R) >> 1;
        if (pos <= M) pst[cur].l = update(pst[old].l, L, M, pos);
        else pst[cur].r = update(pst[old].r, M + 1, R, pos);
    }
    return cur;
}
int kth(int left_root, int right_root, int L, int R, int k) {
    // 两个前缀版本相减，得到原数组区间 [l,r] 的频次
    if (L == R) return L;
    int M = (L + R) >> 1;
    int left_count = pst[pst[right_root].l].sum
                   - pst[pst[left_root].l].sum;
    // 左半边数量足够就向左，否则扣掉左边数量向右找
    if (k <= left_count)
        return kth(pst[left_root].l, pst[right_root].l, L, M, k);
    return kth(pst[left_root].r, pst[right_root].r, M + 1, R,
               k - left_count);
}
\end{lstlisting}

\begin{card}
\textbf{完整样例}\\
\begin{lstlisting}
5 2
5 1 4 4 2
2 5 2
2 5 3
\end{lstlisting}
输出：
\begin{lstlisting}
2
4
\end{lstlisting}
查询前要把离散化排名映射回原值；`k` 不在区间长度内时按题面处理无解。
\end{card}

\section{重链剖分：树上路径和}

\cardmeta{现场常用}{树上路径修改/查询、子树操作}{预处理 \(O(n)\)，单次 \(O(\log^2 n)\)}{DFS 序连续覆盖子树；路径拆成若干重链}

\begin{lstlisting}[style=cpp]
vector<vector<int>> tree;
vector<int> parent, depth, sz, heavy, top, dfn, rev;
int timer;

void dfs1(int u, int p) {
    // 第一次 DFS 求子树大小，并找每个点最大的重儿子
    parent[u] = p; sz[u] = 1; heavy[u] = -1;
    int best = 0;
    for (int v : tree[u]) if (v != p) {
        depth[v] = depth[u] + 1;
        dfs1(v, u); sz[u] += sz[v];
        if (sz[v] > best) best = sz[v], heavy[u] = v;
    }
}
void dfs2(int u, int tp) {
    // 第二次 DFS 给重链连续编号；top[u] 是当前重链顶端
    top[u] = tp; dfn[u] = ++timer; rev[timer] = u;
    if (heavy[u] != -1) dfs2(heavy[u], tp);
    for (int v : tree[u])
        if (v != parent[u] && v != heavy[u]) dfs2(v, v);
}
int path_sum(int u, int v) {
    int ans = 0;
    while (top[u] != top[v]) {
        // 不在同一条重链时，先查询深度更大的链顶并跳到其父亲
        if (depth[top[u]] < depth[top[v]]) swap(u, v);
        ans += seg.query(1, 1, n, dfn[top[u]], dfn[u]);
        u = parent[top[u]];
    }
    if (depth[u] > depth[v]) swap(u, v);
    return ans + seg.query(1, 1, n, dfn[u], dfn[v]);
}
\end{lstlisting}

\begin{warnbox}
这段路径函数依赖基础章节的线段树对象 `seg`。接入时先把点权按 `dfn` 顺序建入线段树；边权问题要把边权挂到更深的端点，并在查询时排除 LCA。
\end{warnbox}

\begin{card}
\textbf{完整样例}\\
树边：`1-2, 1-3, 3-4, 3-5`，点权全为 1。查询路径 `(2,5)` 的和为 `3`；把点 3 加 4 后，再查询输出 `7`。
\end{card}

\section{树上启发式合并 DSU on Tree}

\cardmeta{高风险模板}{每个子树颜色频次、众数和}{通常 \(O(n\log n)\)}{先求重儿子；轻子树统计后必须清空}

\begin{lstlisting}[style=cpp]
vector<vector<int>> tree;
vector<int> color, sz, heavy, ans, cnt;
int big_child, best_freq, best_sum;

void dfs_size(int u, int p) {
    sz[u] = 1; heavy[u] = -1;
    int best = 0;
    for (int v : tree[u]) if (v != p) {
        dfs_size(v, u); sz[u] += sz[v];
        if (sz[v] > best) best = sz[v], heavy[u] = v;
    }
}
void add_subtree(int u, int p, int delta) {
    // delta=+1 加入子树，delta=-1 清空子树
    cnt[color[u]] += delta;
    if (cnt[color[u]] > best_freq ||
        (cnt[color[u]] == best_freq && color[u] < best_sum))
        best_freq = cnt[color[u]], best_sum = color[u];
    for (int v : tree[u])
        if (v != p && v != big_child) add_subtree(v, u, delta);
}
void dfs_sack(int u, int p, bool keep) {
    // 轻儿子处理后清空，重儿子保留，保证总复杂度
    for (int v : tree[u])
        if (v != p && v != heavy[u]) dfs_sack(v, u, false);
    if (heavy[u] != -1) dfs_sack(heavy[u], u, true), big_child = heavy[u];
    add_subtree(u, p, 1);
    big_child = -1;
    ans[u] = best_sum;
    if (!keep) {
        add_subtree(u, p, -1);
        best_freq = best_sum = 0;
    }
}
\end{lstlisting}

\begin{warnbox}
真实题目若要求“最大频率颜色之和”，答案维护器要改成维护所有达到最大频率的颜色和；上面只展示“频率最高且取最小颜色”的示意接口，复制前必须按题意修改。
\end{warnbox}

\begin{card}
\textbf{完整样例}\\
树边：`1-2,1-3,3-4,3-5`，颜色为 `1,2,2,3,2`。\\
每个子树中出现次数最多的颜色（取最小编号）输出：`1 2 2 3 2`。
\end{card}

\section{单调斜率优化 CHT}

\cardmeta{高风险模板}{DP 转移可化为直线、斜率和查询单调}{均摊 \(O(1)\)}{只有在斜率插入顺序和查询顺序满足单调性时使用}

\begin{lstlisting}[style=cpp]
struct Line { int k, b; int value(int x) const { return k * x + b; } };
deque<Line> hull;
bool bad(Line a, Line b, Line c) {
    // b 在几何上被 a、c 遮挡时永远不会成为最优线
    return (i128)(b.b - a.b) * (b.k - c.k)
         >= (i128)(c.b - b.b) * (a.k - b.k);
}
void add_line(Line x) {
    while (hull.size() >= 2 &&
           bad(hull[hull.size()-2], hull.back(), x)) hull.pop_back();
    hull.push_back(x);
}
int query(int x) {
    // 查询点 x 单调递增时，队首劣线可以永久弹出
    while (hull.size() >= 2 &&
           hull[0].value(x) >= hull[1].value(x)) hull.pop_front();
    return hull.front().value(x);
}
\end{lstlisting}

\begin{warnbox}
如果斜率或查询点不单调，不能直接使用这份 deque 版 CHT，应改用李超线段树；本重构不收录李超树的未验证变体。
\end{warnbox}

\begin{card}
\textbf{完整样例}\\
依次加入直线 \(y=1x+0\)、\(y=2x-1\)，按 \(x=0,1,2\) 查询最大值，输出 `0 1 3`。
如果题目的斜率插入顺序或查询顺序不单调，不能套这份代码。
\end{card}
